\documentclass{article}
\usepackage[utf8]{inputenc}
\usepackage[T1]{fontenc}
\usepackage{geometry}
\geometry{
    a4paper,
    left=2.5cm,
    right=2.5cm,
    top=2.5cm,
    bottom=2.5cm
}
\usepackage{graphicx}
\usepackage{float}
\usepackage{hyperref}
\usepackage[english]{babel}
\usepackage{amsmath}
\usepackage{amssymb}
\usepackage{array}
\usepackage{booktabs}
\usepackage{caption}
\usepackage{subcaption}
\usepackage{enumitem}
\usepackage{fancyhdr}
\usepackage{titlesec}
\usepackage{url}
\usepackage{textcomp}
\usepackage{xcolor}
\usepackage{listings}
\usepackage{multirow}
\usepackage{longtable}
\usepackage{tabularx}
\usepackage{times}

% Set up page style
\pagestyle{fancy}
\fancyhf{}
\rfoot{\thepage}
\renewcommand{\headrulewidth}{0pt}

% Title formatting
\titleformat{\chapter}[display]{\normalfont\bfseries}{\chaptertitle\thechapter}{0pt}{\Huge}
\titlespacing*{\chapter}{0pt}{50pt}{40pt}

% Document info
\title{Template Document}
\author{Academic Support Agent}
\date{\today}

\begin{document}
\maketitle

\title{Anota es\ da\ aula}
%% {\large \today}\\[5mm] %% se quiser colocar data

\section*{Template para anota es de aula}

\subsection*{Aula 1}


%% para citar no texto, use \textcite{citacao} para ficar no formato Fulano (Ano), ou use \cite{citacao} para citar  no formato (FULANO, Ano)
Segundo \textcite{citacao-exemplo}, \lipsum[66] \cite{citacao-exemplo}

\par
Para equa es, listas, etc., basta usar os environments padr o do LaTeX.

\begin{equation}
a^2 = b^2 + c^2
\end{equation}
onde $a$, $b$ e $c$ s o lados de um tri ngulo ret ngulo........

\subsubsection*{Subse o da aula 1}

\lipsum[66]

\myparagraph{Par grafo da subse o da aula 1}
\lipsum[66]
\par A \autoref{fig:exemplo}   uma imagem......

\begin{figure}[!htb]
    \centering
    \caption{Figura exemplo}
    \includegraphics[width=0.5\textwidth]{logo_utf.png}
    \caption*{Fonte: Autoria pr pria.} %% isso   levemente uma gambiarra, mas funciona para o prop sito desse template.
    \label{fig:exemplo}
\end{figure}

\mysubparagraph{Subpar grafo da subse o da aula 1}
\lipsum[66]

%%% as refer ncias devem estar em formato bibTeX no arquivo referencias.bib
\printbibliography
\end{document}